\section{参加人}
\begin{enumerate}
  \item 能不能做被告。
  \item 直接起诉被告的选择问题。
        \begin{enumerate}
          \item 原则\bs{谁行为谁被告（即使越权）}。
          \item \bs{临时机构、经批准、被、委、派、和、垂直总队、共同、开发、房屋}。
          \item 临时机构。设立者、组建者。
          \item 经批准行为诉讼看名义（盖章）。\al{经\xhx{市公安局}批准……县公安局罚款王某 200 元}。常考，背例子。
          \item 被授权方、委托机关、派出机关 \yjt 没了找继承者。
          \item 垂直领导找上级。\al{国安、税务}。
          \item 总队 \dy 支队 \dy 大队（超纲，了解）\ydy 基层政府的局。\al{刑侦大队、消防支队}。
          \item 共同行为告所有。\al{工商 + 环保}。
                \begin{enumerate}
                  \item \bs{告漏了 \yjt 追加为第三人}。和民诉不一样。
                  \item \bs{告错 \yjt 告知变更} \yjt 不愿变更 \yjt 裁驳。\bs{想变就变，不变就驳}。
                  \item \zd{假共同 \yjt 去假留真}。\al{工商局和消协作的处罚}。
                \end{enumerate}
          \item 开发区。是/不是\zd{国务院 \hz 省政府}批准设立的 \yjt \ydy 区政府/临时机构（\zd{是否有授权}）。
          \item 房屋征收补偿行为。房屋征收部门是政府下的一个部门（\al{建设局/建委/规划与建设局/土地与建设局）}。
          \item 内设/派出机构有授权。\bs{小祸自己扛，大祸爹来挡}。考\zd{派出所}。罚了 800/拘留 \yjt 幅度/种类越权。
          \item \al{选项同时出现县政府，县环保局，选职能部门}。原则是尽量往下压案子。新增。
                \begin{enumerate}
                  \item 指令/责\zd{令} = 命令 \yjt 县政府。
                  \item \zd{责}成 = 交办 \yjt \bs{有文书看文书 \yjt 没文书看行动}（大概率是职能部门）。
                  \item \zd{指导}（温柔的建议）\yjt 职能部门。
                \end{enumerate}
        \end{enumerate}
        \begin{kk}{先复议后起诉被告的选择问题}
          \begin{enumerate}
            \item \bs{复议不作为择一告}。复议不作为（\zd{没有实质审理过}）\yjt 复议机关。原行为 \yjt 原机关。
                  \al{区政府征收甲房屋，甲提起复议，市政府不予受理，甲提起诉讼}。
            \item \bs{复议改变单独告}。改变\zd{实体权利义务、结果}。\al{动证据、改法条、纠程序 \cuo}。
            \item \bs{复议维持共同告}。
                  \al{甲申请工伤认定，区人社局《不予受理决定书》，甲找区政府\xhx{复议}，给了一个《不予认定工伤复议决定书》，甲以区人社局为被告提起诉讼}。
                  \begin{enumerate}
                    \item \zd{审过 + 结果没动}。\al{不予受理 \yjt 本质就是不予工伤认定}。\bs{注意变相维持、判驳}。
                    \item \zd{告漏}通知原告变更或追加。不同意追加 \yjt \zd{强行追加共同被告}（\al{追加为第三人 \cuo}）。
                    \item 复议维持本质上是\zd{重复处理}的一种。题目一定会出现“复议”两字。
                  \end{enumerate}
            \item \bs{多行为、多请求、多因素共同告}。一个复议决定里既有复议改变，又有复议维持。
          \end{enumerate}
          \tcblower
          \zd{排除法。审过了排除不作为，结果没动排除改变}。
        \end{kk}
  \item 原告。\bs{认为行政行为侵犯其合法权益而依法向法院起诉的公民、法人、其他组织}。
        \begin{enumerate}
          \item 认为 \yjt 只看主观\zd{不看客观}。
          \item 其 \yjt 自己。\zd{不能为了他人、公共利益}（个人\zd{不能}提起公益诉讼）。
          \item 侵犯 \yjt \zd{必然影响/限制/侵犯}。
          \item 合法 \yjt \al{毒品 \cuo}。
          \item 其他组织 \yjt \al{政府机关做原告 \cuo}。
                \begin{kk}{判断标准}
                  \begin{enumerate}
                    \item \zd{相对人}必然具有原告资格，必然有利害关系。
                    \item 相关人仅限\bs{侵权/物权关系、公平竟争对手、亲属关系 \cuo（死后继受有）、合同能民诉无行诉}。
                          \al{乙营业执照被吊销，致使和丙之间的买卖合同无法履行，丙有权提起行政诉讼/第三人资格 \cuo}。
                    \item 不公平竟争。\al{A 幼儿园，BC 申请，只给了 B，C \dui A \cuo 不保障垄断（A）的既得利益}。
                  \end{enumerate}
                \end{kk}
        \end{enumerate}
  \item 组织。\bs{股份、联合、死企业；非营、业主、合、个体}。
        \begin{enumerate}
          \item 股份制。\zd{3 加 1}。董事会、股东大会、股东会/股东代表大会 + 法代。\al{股东 \cuo 和股权占比无关}。
          \item 联营、合资、合作。法代、\zd{投资人自己名义}。
          \item 死了的企业。\zd{原/新企业、法代}。\zd{死了 = 注销、合并、强令兼并、出售、分立、改变隶属关系}。
          \item 非营利/营利法人。\zd{投资人自己名义}。
          \item 业主委员会。业主。
          \item 合伙和个体工商案件。
        \end{enumerate}
  \item 共同诉讼。基本同民诉。
        \begin{enumerate}
          \item 告漏了。
                \begin{enumerate}
                  \item 必要的/普通的 \yjt 同一\zd{个/类}行政行为 \yjt \zd{应当/可以}通知。
                  \item 追加共同被告。\al{复议维持}。
                  \item 追加共同原告（基于同一个法律事实关系，利益一致）。
                        \begin{enumerate}
                          \item 通知 + 同意 = \zd{共同原告}。
                          \item 通知 + 不同意 + \zd{不弃权}/肯弃权 = \zd{第三人}/不追加（啥也不做）。
                        \end{enumerate}
                \end{enumerate}
          \item 第三人。\zd{利害关系人起诉了是原告，没起诉做第三人}。\zd{行诉里均为有独三}。
                \begin{enumerate}
                  \item 可以转换的含义。不能做第三人/原告 = 不能做原告/第三人。
                  \item 分原告型（民）、\zd{被告型}（行政机关）第三人。经常考，就 4 种，死记。
                        \begin{enumerate}
                          \item \bs{共同、改变、矛盾}。假共同、共同行为告漏、复议改变原机关、多个机关相互矛盾但非被告。
                          \item 不来。缺席审判。
                        \end{enumerate}
                  \item 诉讼地位。\bs{一审、依申请；能举证、能上诉、能参加调解、能申请执行}。
                \end{enumerate}
          \item 直接追加的情形。与\zd{原告利益相关}（判决预决效力，结果上的）的利害关系人。
          \item 代理人。最主要的考点，记\zd{多少个人要选代表}。
                \begin{enumerate}
                  \item 人数。\bs{二五一十，一五得五}。一谐音复议。
                  \item 自己选、法院指定、委托。
                  \item 代表范围。律师、近亲属、员工、工作人员、社会团体推荐的。推荐出来的\zd{公民代理没资格}的。
                  \item 方式。原则授权委托书。例外。\zd{被限制人身自由可以口头}。
                  \item 解除委托。自行通知其他当事人。\al{法院通知 \cuo 法院永远不是主动方}。
                \end{enumerate}
        \end{enumerate}
\end{enumerate}
